\documentclass[10pt]{article}
% Hojas resumen: un A4 por delante y por detrás, dos columnas por cara.
% Uso: \documentclass[10pt]{article} % Hojas resumen: un A4 por delante y por detrás, dos columnas por cara.
% Uso: \documentclass[10pt]{article} % Hojas resumen: un A4 por delante y por detrás, dos columnas por cara.
% Uso: \documentclass[10pt]{article} \input{estilo-resumen.tex} \bloque{analisis|algebra|azar} \cabecera{n}{Título}{Lecciones}
\usepackage[utf8]{inputenc}
\usepackage[T1]{fontenc}
\usepackage[spanish,es-nodecimaldot]{babel}
\usepackage{lmodern}
\usepackage{amsmath,amssymb}
\usepackage{graphicx}
\usepackage[a4paper,left=10mm,right=10mm,top=15mm,bottom=13mm,headheight=11pt,headsep=3mm,footskip=6mm]{geometry}
\usepackage{multicol}
\setlength{\columnsep}{6mm}
\setlength{\columnseprule}{0.3pt}
\usepackage{xcolor}
\usepackage{enumitem}
\setlist{nosep,leftmargin=1.1em,itemsep=1pt}
\usepackage{array}
\usepackage{needspace}
\usepackage{pgfplots}
\usepgfplotslibrary{fillbetween}
\pgfplotsset{compat=1.18,
  mini/.style={width=4.3cm,height=3.3cm,scale only axis=false,axis lines=middle,xtick=\empty,ytick=\empty,samples=120,
    xlabel={\scriptsize$x$},ylabel={\scriptsize$y$},title style={font=\scriptsize,yshift=-4pt},clip=true,
    every axis plot/.append style={thick,color=acento}}}
\usepackage{ragged2e}
\usepackage[most]{tcolorbox}
\usepackage{fancyhdr,lastpage}
\usepackage[hidelinks,pdfencoding=auto]{hyperref}
\hypersetup{pdflang=es-ES,pdfsubject={Matemáticas Aplicadas a las Ciencias Sociales II}}

\definecolor{analisis}{HTML}{1D4ED8}
\definecolor{algebra}{HTML}{86198F}
\definecolor{azar}{HTML}{9A3412}
\definecolor{aviso}{HTML}{B91C1C}
\colorlet{acento}{analisis}
\newcommand{\bloque}[1]{\colorlet{acento}{#1}}

\setlength{\parindent}{0pt}
\setlength{\parskip}{2pt}
\renewcommand{\arraystretch}{1.2}
\AtBeginDocument{\RaggedRight\fontsize{9.8}{11.9}\selectfont\setlength{\abovedisplayskip}{3pt}\setlength{\belowdisplayskip}{3pt}\setlength{\abovedisplayshortskip}{1pt}\setlength{\belowdisplayshortskip}{1pt}}

\newcommand{\temanum}{}\newcommand{\tematitulo}{}
\pagestyle{fancy}
\fancyhf{}
\fancyhead[L]{\footnotesize\color{acento}\textbf{Hoja resumen}\quad Tema \temanum. \tematitulo}
\fancyhead[R]{\footnotesize Matemáticas Aplicadas a las Ciencias Sociales II --- 2.\textsuperscript{o} Bachillerato}
\fancyfoot[L]{\scriptsize IES Al-Ándalus --- Juan María Gámez Ortiz, con la ayuda de Claude y de Codex --- CC BY-NC 4.0}
\fancyfoot[R]{\scriptsize Cara \thepage\ de \pageref{LastPage}}
\renewcommand{\headrulewidth}{0.4pt}

% Cabecera de la primera cara: número, título y lecciones del tema.
\newcommand{\cabecera}[3]{%
  \renewcommand{\temanum}{#1}\renewcommand{\tematitulo}{#2}%
  \hypersetup{pdftitle={Hoja resumen. Tema #1. #2}}%
  \noindent\begin{tcolorbox}[enhanced,colback=acento!8,colframe=acento,boxrule=0.6pt,arc=2pt,left=4pt,right=4pt,top=2pt,bottom=2pt]
  \begin{minipage}[c]{1.1cm}\includegraphics[width=1cm]{logo-al-andalus.png}\end{minipage}%
  \begin{minipage}[c]{\dimexpr\linewidth-1.2cm}{\large\bfseries\color{acento}Tema #1. #2}\\[1pt]{\scriptsize #3}\end{minipage}
  \end{tcolorbox}\vspace{-2pt}}

% Apartado con filete del color del bloque.
\newcommand{\seccion}[1]{\par\needspace{5\baselineskip}\vspace{3pt}{\color{acento}\bfseries\normalsize #1}\par\vspace{-1pt}{\color{acento}\hrule height 0.5pt}\vspace{2pt}}
\newcommand{\sub}[1]{\par\needspace{3\baselineskip}\vspace{1pt}\textbf{\color{acento}#1}\enspace}

% Fórmula clave, aviso de errores y ejemplo.
\tcbset{before upper=\RaggedRight}
\newtcolorbox{formula}{enhanced,colback=acento!6,colframe=acento!70,boxrule=0.5pt,arc=1.5pt,left=3pt,right=3pt,top=0pt,bottom=0pt,before skip=2pt,after skip=2pt}
\newtcolorbox{ojo}{enhanced,colback=aviso!5,colframe=aviso!60,boxrule=0.5pt,arc=1.5pt,left=3pt,right=3pt,top=1pt,bottom=1pt,before skip=3pt,after skip=2pt,
  fonttitle=\bfseries\footnotesize,coltitle=aviso,colbacktitle=aviso!5,title={Ojo},attach title to upper={\enspace}}
\newtcolorbox{ejemplo}{enhanced,colback=white,colframe=acento!45,boxrule=0.4pt,arc=1.5pt,left=3pt,right=3pt,top=1pt,bottom=1pt,before skip=3pt,after skip=2pt,
  fonttitle=\bfseries\footnotesize,coltitle=acento,colbacktitle=white,title={Ejemplo},attach title to upper={\enspace}}
\newcommand{\dx}{\,dx}
 \bloque{analisis|algebra|azar} \cabecera{n}{Título}{Lecciones}
\usepackage[utf8]{inputenc}
\usepackage[T1]{fontenc}
\usepackage[spanish,es-nodecimaldot]{babel}
\usepackage{lmodern}
\usepackage{amsmath,amssymb}
\usepackage{graphicx}
\usepackage[a4paper,left=10mm,right=10mm,top=15mm,bottom=13mm,headheight=11pt,headsep=3mm,footskip=6mm]{geometry}
\usepackage{multicol}
\setlength{\columnsep}{6mm}
\setlength{\columnseprule}{0.3pt}
\usepackage{xcolor}
\usepackage{enumitem}
\setlist{nosep,leftmargin=1.1em,itemsep=1pt}
\usepackage{array}
\usepackage{needspace}
\usepackage{pgfplots}
\usepgfplotslibrary{fillbetween}
\pgfplotsset{compat=1.18,
  mini/.style={width=4.3cm,height=3.3cm,scale only axis=false,axis lines=middle,xtick=\empty,ytick=\empty,samples=120,
    xlabel={\scriptsize$x$},ylabel={\scriptsize$y$},title style={font=\scriptsize,yshift=-4pt},clip=true,
    every axis plot/.append style={thick,color=acento}}}
\usepackage{ragged2e}
\usepackage[most]{tcolorbox}
\usepackage{fancyhdr,lastpage}
\usepackage[hidelinks,pdfencoding=auto]{hyperref}
\hypersetup{pdflang=es-ES,pdfsubject={Matemáticas Aplicadas a las Ciencias Sociales II}}

\definecolor{analisis}{HTML}{1D4ED8}
\definecolor{algebra}{HTML}{86198F}
\definecolor{azar}{HTML}{9A3412}
\definecolor{aviso}{HTML}{B91C1C}
\colorlet{acento}{analisis}
\newcommand{\bloque}[1]{\colorlet{acento}{#1}}

\setlength{\parindent}{0pt}
\setlength{\parskip}{2pt}
\renewcommand{\arraystretch}{1.2}
\AtBeginDocument{\RaggedRight\fontsize{9.8}{11.9}\selectfont\setlength{\abovedisplayskip}{3pt}\setlength{\belowdisplayskip}{3pt}\setlength{\abovedisplayshortskip}{1pt}\setlength{\belowdisplayshortskip}{1pt}}

\newcommand{\temanum}{}\newcommand{\tematitulo}{}
\pagestyle{fancy}
\fancyhf{}
\fancyhead[L]{\footnotesize\color{acento}\textbf{Hoja resumen}\quad Tema \temanum. \tematitulo}
\fancyhead[R]{\footnotesize Matemáticas Aplicadas a las Ciencias Sociales II --- 2.\textsuperscript{o} Bachillerato}
\fancyfoot[L]{\scriptsize IES Al-Ándalus --- Juan María Gámez Ortiz, con la ayuda de Claude y de Codex --- CC BY-NC 4.0}
\fancyfoot[R]{\scriptsize Cara \thepage\ de \pageref{LastPage}}
\renewcommand{\headrulewidth}{0.4pt}

% Cabecera de la primera cara: número, título y lecciones del tema.
\newcommand{\cabecera}[3]{%
  \renewcommand{\temanum}{#1}\renewcommand{\tematitulo}{#2}%
  \hypersetup{pdftitle={Hoja resumen. Tema #1. #2}}%
  \noindent\begin{tcolorbox}[enhanced,colback=acento!8,colframe=acento,boxrule=0.6pt,arc=2pt,left=4pt,right=4pt,top=2pt,bottom=2pt]
  \begin{minipage}[c]{1.1cm}\includegraphics[width=1cm]{logo-al-andalus.png}\end{minipage}%
  \begin{minipage}[c]{\dimexpr\linewidth-1.2cm}{\large\bfseries\color{acento}Tema #1. #2}\\[1pt]{\scriptsize #3}\end{minipage}
  \end{tcolorbox}\vspace{-2pt}}

% Apartado con filete del color del bloque.
\newcommand{\seccion}[1]{\par\needspace{5\baselineskip}\vspace{3pt}{\color{acento}\bfseries\normalsize #1}\par\vspace{-1pt}{\color{acento}\hrule height 0.5pt}\vspace{2pt}}
\newcommand{\sub}[1]{\par\needspace{3\baselineskip}\vspace{1pt}\textbf{\color{acento}#1}\enspace}

% Fórmula clave, aviso de errores y ejemplo.
\tcbset{before upper=\RaggedRight}
\newtcolorbox{formula}{enhanced,colback=acento!6,colframe=acento!70,boxrule=0.5pt,arc=1.5pt,left=3pt,right=3pt,top=0pt,bottom=0pt,before skip=2pt,after skip=2pt}
\newtcolorbox{ojo}{enhanced,colback=aviso!5,colframe=aviso!60,boxrule=0.5pt,arc=1.5pt,left=3pt,right=3pt,top=1pt,bottom=1pt,before skip=3pt,after skip=2pt,
  fonttitle=\bfseries\footnotesize,coltitle=aviso,colbacktitle=aviso!5,title={Ojo},attach title to upper={\enspace}}
\newtcolorbox{ejemplo}{enhanced,colback=white,colframe=acento!45,boxrule=0.4pt,arc=1.5pt,left=3pt,right=3pt,top=1pt,bottom=1pt,before skip=3pt,after skip=2pt,
  fonttitle=\bfseries\footnotesize,coltitle=acento,colbacktitle=white,title={Ejemplo},attach title to upper={\enspace}}
\newcommand{\dx}{\,dx}
 \bloque{analisis|algebra|azar} \cabecera{n}{Título}{Lecciones}
\usepackage[utf8]{inputenc}
\usepackage[T1]{fontenc}
\usepackage[spanish,es-nodecimaldot]{babel}
\usepackage{lmodern}
\usepackage{amsmath,amssymb}
\usepackage{graphicx}
\usepackage[a4paper,left=10mm,right=10mm,top=15mm,bottom=13mm,headheight=11pt,headsep=3mm,footskip=6mm]{geometry}
\usepackage{multicol}
\setlength{\columnsep}{6mm}
\setlength{\columnseprule}{0.3pt}
\usepackage{xcolor}
\usepackage{enumitem}
\setlist{nosep,leftmargin=1.1em,itemsep=1pt}
\usepackage{array}
\usepackage{needspace}
\usepackage{pgfplots}
\usepgfplotslibrary{fillbetween}
\pgfplotsset{compat=1.18,
  mini/.style={width=4.3cm,height=3.3cm,scale only axis=false,axis lines=middle,xtick=\empty,ytick=\empty,samples=120,
    xlabel={\scriptsize$x$},ylabel={\scriptsize$y$},title style={font=\scriptsize,yshift=-4pt},clip=true,
    every axis plot/.append style={thick,color=acento}}}
\usepackage{ragged2e}
\usepackage[most]{tcolorbox}
\usepackage{fancyhdr,lastpage}
\usepackage[hidelinks,pdfencoding=auto]{hyperref}
\hypersetup{pdflang=es-ES,pdfsubject={Matemáticas Aplicadas a las Ciencias Sociales II}}

\definecolor{analisis}{HTML}{1D4ED8}
\definecolor{algebra}{HTML}{86198F}
\definecolor{azar}{HTML}{9A3412}
\definecolor{aviso}{HTML}{B91C1C}
\colorlet{acento}{analisis}
\newcommand{\bloque}[1]{\colorlet{acento}{#1}}

\setlength{\parindent}{0pt}
\setlength{\parskip}{2pt}
\renewcommand{\arraystretch}{1.2}
\AtBeginDocument{\RaggedRight\fontsize{9.8}{11.9}\selectfont\setlength{\abovedisplayskip}{3pt}\setlength{\belowdisplayskip}{3pt}\setlength{\abovedisplayshortskip}{1pt}\setlength{\belowdisplayshortskip}{1pt}}

\newcommand{\temanum}{}\newcommand{\tematitulo}{}
\pagestyle{fancy}
\fancyhf{}
\fancyhead[L]{\footnotesize\color{acento}\textbf{Hoja resumen}\quad Tema \temanum. \tematitulo}
\fancyhead[R]{\footnotesize Matemáticas Aplicadas a las Ciencias Sociales II --- 2.\textsuperscript{o} Bachillerato}
\fancyfoot[L]{\scriptsize IES Al-Ándalus --- Juan María Gámez Ortiz, con la ayuda de Claude y de Codex --- CC BY-NC 4.0}
\fancyfoot[R]{\scriptsize Cara \thepage\ de \pageref{LastPage}}
\renewcommand{\headrulewidth}{0.4pt}

% Cabecera de la primera cara: número, título y lecciones del tema.
\newcommand{\cabecera}[3]{%
  \renewcommand{\temanum}{#1}\renewcommand{\tematitulo}{#2}%
  \hypersetup{pdftitle={Hoja resumen. Tema #1. #2}}%
  \noindent\begin{tcolorbox}[enhanced,colback=acento!8,colframe=acento,boxrule=0.6pt,arc=2pt,left=4pt,right=4pt,top=2pt,bottom=2pt]
  \begin{minipage}[c]{1.1cm}\includegraphics[width=1cm]{logo-al-andalus.png}\end{minipage}%
  \begin{minipage}[c]{\dimexpr\linewidth-1.2cm}{\large\bfseries\color{acento}Tema #1. #2}\\[1pt]{\scriptsize #3}\end{minipage}
  \end{tcolorbox}\vspace{-2pt}}

% Apartado con filete del color del bloque.
\newcommand{\seccion}[1]{\par\needspace{5\baselineskip}\vspace{3pt}{\color{acento}\bfseries\normalsize #1}\par\vspace{-1pt}{\color{acento}\hrule height 0.5pt}\vspace{2pt}}
\newcommand{\sub}[1]{\par\needspace{3\baselineskip}\vspace{1pt}\textbf{\color{acento}#1}\enspace}

% Fórmula clave, aviso de errores y ejemplo.
\tcbset{before upper=\RaggedRight}
\newtcolorbox{formula}{enhanced,colback=acento!6,colframe=acento!70,boxrule=0.5pt,arc=1.5pt,left=3pt,right=3pt,top=0pt,bottom=0pt,before skip=2pt,after skip=2pt}
\newtcolorbox{ojo}{enhanced,colback=aviso!5,colframe=aviso!60,boxrule=0.5pt,arc=1.5pt,left=3pt,right=3pt,top=1pt,bottom=1pt,before skip=3pt,after skip=2pt,
  fonttitle=\bfseries\footnotesize,coltitle=aviso,colbacktitle=aviso!5,title={Ojo},attach title to upper={\enspace}}
\newtcolorbox{ejemplo}{enhanced,colback=white,colframe=acento!45,boxrule=0.4pt,arc=1.5pt,left=3pt,right=3pt,top=1pt,bottom=1pt,before skip=3pt,after skip=2pt,
  fonttitle=\bfseries\footnotesize,coltitle=acento,colbacktitle=white,title={Ejemplo},attach title to upper={\enspace}}
\newcommand{\dx}{\,dx}

\bloque{azar}
\begin{document}
\cabecera{8}{Distribuciones de probabilidad}{Lecciones 8.1--8.4: variable aleatoria discreta y continua, distribución binomial, distribución normal y aproximación de la binomial por la normal.}
\begin{multicols}{2}

\seccion{1. Variable aleatoria}
Una \textbf{variable aleatoria} \(X\) asigna un número real a cada resultado del experimento: \(X:\Omega\to\mathbb R\).
\begin{itemize}
  \item \textbf{Discreta}: valores que se pueden enumerar (finitos o numerables), aunque no sean enteros: un premio de 0,50 o 1,50 euros.
  \item \textbf{Continua}: se describe con una densidad y probabilidades de intervalos (tiempos, pesos).
\end{itemize}
\sub{Distribución discreta finita.}
\begin{formula}
\[P(X=x_i)=p_i,\qquad 0\le p_i\le1,\qquad \sum p_i=1\]
\[\mu=\sum x_ip_i\]
\[\sigma^2=\sum x_i^2p_i-\mu^2,\qquad \sigma=\sqrt{\sigma^2}\]
\end{formula}
\begin{ejemplo}
Dado equilibrado: \(\mu=\frac{21}6=3{,}5\); \(\sum x_i^2p_i=\frac{91}6\); \(\sigma^2=\frac{91}6-\frac{49}4=\frac{35}{12}\approx2{,}92\); \(\sigma\approx1{,}71\). La media no tiene por qué ser un valor posible.
\end{ejemplo}
\sub{Densidad} (continua).
\begin{formula}
\[f(x)\ge0,\qquad \int_{-\infty}^{+\infty}f(x)\,dx=1\]
\[P(a\le X\le b)=\int_a^bf(x)\,dx\]
\end{formula}
La probabilidad es un \textbf{área}: \(P(X=a)=0\), y da igual incluir o no los extremos. La altura \(f(x)\) no es una probabilidad (puede ser mayor que 1).
\sub{Acumulada:} \(F(x)=P(X\le x)\), no decreciente, de 0 a 1. Discreta: suma de masas; continua: área hasta \(x\).

\seccion{2. Distribución binomial}
Condiciones: \(n\) ensayos \textbf{fijos}, dos resultados (éxito/fracaso), \textbf{independientes} y con la \textbf{misma} probabilidad de éxito \(p\). \(X\) = número de éxitos.
\begin{formula}
\[X\sim B(n,p):\qquad P(X=k)=\binom nk p^kq^{n-k},\quad q=1-p\]
\[\mu=np,\qquad \sigma=\sqrt{npq}\]
\end{formula}
\(\binom nk\) cuenta las posiciones posibles de los \(k\) éxitos (nCr en la calculadora).
\begin{ejemplo}
Encesta el 80\,\%; lanza 5 veces. \(X\sim B(5;\,0{,}8)\).
\[P(X=3)=\binom53(0{,}8)^3(0{,}2)^2=10\cdot0{,}512\cdot0{,}04=0{,}2048\]
\(\mu=4\) aciertos, \(\sigma=\sqrt{0{,}8}\approx0{,}89\).
\end{ejemplo}
\sub{Traducir con la acumulada \(F(k)=P(X\le k)\).}
\begin{tabular}{@{}lll@{}}
Como máximo \(k\) & \(X\le k\) & \(F(k)\)\\
Menos de \(k\) & \(X<k\) & \(F(k-1)\)\\
Al menos \(k\) & \(X\ge k\) & \(1-F(k-1)\)\\
Más de \(k\) & \(X>k\) & \(1-F(k)\)\\
Entre \(a\) y \(b\) & \(a\le X\le b\) & \(F(b)-F(a-1)\)
\end{tabular}

\begin{ejemplo}
\(B(5;\,0{,}8)\): \(P(X\ge3)=1-P(X\le2)=1-0{,}05792=0{,}94208\). Al menos uno: \(1-0{,}2^5=0{,}99968\).
\end{ejemplo}

\seccion{3. Distribución normal}
\(X\sim N(\mu,\sigma)\): campana simétrica respecto a \(\mu\); \(\sigma\) (\textbf{desviación típica}, el segundo parámetro) fija su anchura. Puntos de inflexión en \(\mu\pm\sigma\).
\begin{center}
\begin{tikzpicture}
\begin{axis}[mini,width=6.8cm,height=3.4cm,xmin=-3.6,xmax=3.6,ymin=0,ymax=0.45,xtick={-2,-1,0,1,2},xticklabels={\(\mu-2\sigma\),\(\mu-\sigma\),\(\mu\),\(\mu+\sigma\),\(\mu+2\sigma\)},tick label style={font=\tiny},xlabel={},ylabel={},axis y line=none]
\addplot[name path=c,domain=-3.6:3.6]{exp(-x^2/2)/sqrt(2*pi)};
\path[name path=ax] (axis cs:-3.6,0) -- (axis cs:3.6,0);
\addplot[fill=acento!30,draw=none] fill between[of=c and ax,soft clip={domain=-1:1}];
\addplot[fill=acento!12,draw=none] fill between[of=c and ax,soft clip={domain=-2:-1}];
\addplot[fill=acento!12,draw=none] fill between[of=c and ax,soft clip={domain=1:2}];
\end{axis}
\end{tikzpicture}
\end{center}
\textbf{Regla empírica:} \(\mu\pm\sigma\): 68,27\,\%; \(\mu\pm2\sigma\): 95,45\,\%; \(\mu\pm3\sigma\): 99,73\,\%.
\begin{formula}
\textbf{Tipificación:}
\[Z=\frac{X-\mu}{\sigma}\sim N(0,1),\qquad \Phi(z)=P(Z\le z)\]
\[P(a\le X\le b)=\Phi\Bigl(\frac{b-\mu}\sigma\Bigr)-\Phi\Bigl(\frac{a-\mu}\sigma\Bigr)\]
\end{formula}
\begin{tabular}{@{}ll@{}}
Cola izquierda & \(P(Z\le z)=\Phi(z)\)\\
Cola derecha & \(P(Z>z)=1-\Phi(z)\)\\
Valores negativos & \(\Phi(-z)=1-\Phi(z)\)\\
Intervalo & \(\Phi(b)-\Phi(a)\)
\end{tabular}

\sub{Extracto de la tabla de \(\Phi\).}
\[\begin{array}{c|cccc}
z&0&1&1{,}28&1{,}64\\\hline
\Phi(z)&0{,}5&0{,}8413&0{,}8997&0{,}9495
\end{array}\]
\[\begin{array}{c|cccc}
z&1{,}65&1{,}96&2&2{,}58\\\hline
\Phi(z)&0{,}9505&0{,}975&0{,}9772&0{,}9951
\end{array}\]
\begin{ejemplo}
Puntuaciones \(X\sim N(100,15)\).
\(P(X<85)=\Phi(-1)=1-0{,}8413=0{,}1587\).
\(P(X>130)=1-\Phi(2)=0{,}0228\).
\(P(85\le X\le130)=0{,}9772-0{,}1587=0{,}8185\).
\end{ejemplo}
\sub{Problema inverso (percentil).} Se busca \(z_r\) con \(\Phi(z_r)=r\) y se deshace la tipificación:
\begin{formula}
\[x=\mu+\sigma z_r\]
\end{formula}
\begin{ejemplo}
Percentil 90 de \(N(100,15)\): \(\Phi(1{,}28)\approx0{,}90\Rightarrow x\approx100+15\cdot1{,}28=119{,}2\). Si se da una cola derecha \(q\), se busca \(\Phi(z)=1-q\).
\end{ejemplo}

\seccion{4. Aproximación de la binomial por la normal}
\begin{formula}
Si \(np\ge5\) y \(n(1-p)\ge5\):
\[X\sim\operatorname{Bin}(n,p)\ \approx\ Y\sim N\bigl(np,\ \sqrt{np(1-p)}\bigr)\]
\end{formula}
\sub{Corrección de continuidad:} cada entero \(k\) ocupa la celda \([k-0{,}5,\ k+0{,}5]\).
\begin{tabular}{@{}ll@{}}
\(P(X\le k)\) & \(\approx P(Y\le k+0{,}5)\)\\
\(P(X<k)\) & \(\approx P(Y\le k-0{,}5)\)\\
\(P(X\ge k)\) & \(\approx P(Y\ge k-0{,}5)\)\\
\(P(X>k)\) & \(\approx P(Y\ge k+0{,}5)\)\\
\(P(A\le X\le B)\) & \(\approx P(A-0{,}5\le Y\le B+0{,}5)\)\\
\(P(X=k)\) & \(\approx P(k-0{,}5\le Y\le k+0{,}5)\)
\end{tabular}

\begin{ejemplo}
\(\operatorname{Bin}(100;\,0{,}5)\): \(np=50\), \(\sigma=5\), \(Y\sim N(50,5)\).
\[P(X\le55)\approx\Phi\Bigl(\frac{55{,}5-50}5\Bigr)=\Phi(1{,}1)\approx0{,}8643\]
(valor binomial exacto: 0,8644).
\[P(X\ge60)\approx1-\Phi(1{,}9)\approx0{,}0287\quad(\text{exacto } 0{,}0284)\]
\[P(45\le X\le55)\approx\Phi(1{,}1)-\Phi(-1{,}1)\approx0{,}7287\]
\[P(X=50)\approx\Phi(0{,}1)-\Phi(-0{,}1)\approx0{,}0797\quad(\text{exacto } 0{,}0796)\]
\end{ejemplo}

\seccion{5. Ejercicio binomial completo}
El 10\,\% de las piezas de un lote son defectuosas. Se revisan 10 al azar (independencia). \(X\sim B(10;\,0{,}1)\).
\begin{center}
\begin{tikzpicture}
\begin{axis}[mini,width=6.4cm,height=3.2cm,ybar,bar width=7pt,xmin=-0.7,xmax=6.7,ymin=0,ymax=0.45,axis x line=bottom,axis y line=none,xtick={0,...,6},tick label style={font=\tiny},xlabel={},ylabel={},title={\(P(X=k)\) en \(B(10;\,0{,}1)\)}]
\addplot[fill=acento!35,draw=acento] coordinates {(0,0.3487) (1,0.3874) (2,0.1937) (3,0.0574) (4,0.0112) (5,0.0015) (6,0.0001)};
\end{axis}
\end{tikzpicture}
\end{center}
\begin{itemize}
  \item Ninguna defectuosa: \(P(X=0)=0{,}9^{10}\approx0{,}3487\).
  \item Como máximo una: \(0{,}3487+10\cdot0{,}1\cdot0{,}9^9\approx0{,}3487+0{,}3874=0{,}7361\).
  \item Al menos dos: \(1-P(X\le1)\approx0{,}2639\).
  \item Media: \(\mu=10\cdot0{,}1=1\) pieza; \(\sigma=\sqrt{0{,}9}\approx0{,}95\).
\end{itemize}

\seccion{6. Ejercicio normal completo}
Los paquetes pesan \(X\sim N(250,10)\) gramos.
\begin{enumerate}
  \item Menos de 240 g: \(z=\frac{240-250}{10}=-1\), \(P=\Phi(-1)=0{,}1587\): el 15,87\,\%.
  \item Entre 240 y 270 g: \(\Phi(2)-\Phi(-1)=0{,}9772-0{,}1587=0{,}8185\).
  \item Peso que solo supera el 5\,\% más pesado: \(\Phi(z)=0{,}95\Rightarrow z\approx1{,}645\), \(x=250+10\cdot1{,}645\approx266{,}5\) g.
\end{enumerate}
\textbf{Pasos:} 1) define \(X\) y sus parámetros; 2) escribe la probabilidad pedida; 3) tipifica; 4) dibuja la campana y decide la cola; 5) usa la tabla y la simetría; 6) responde en el contexto.

\seccion{7. Qué distribución usar}
\begin{tabular}{@{}p{3.6cm}p{4.2cm}@{}}
\textbf{Situación} & \textbf{Modelo}\\ \hline
Contar éxitos en \(n\) ensayos iguales e independientes & Binomial \(B(n,p)\)\\
Medida continua con forma de campana & Normal \(N(\mu,\sigma)\)\\
Binomial con \(n\) grande, \(np\) y \(nq\ge5\) & Normal con corrección\\
Lista de valores con probabilidades & Discreta finita: tabla
\end{tabular}

\begin{ojo}
\begin{itemize}
  \item Usar la binomial sin independencia (extracciones sin reemplazamiento).
  \item Olvidar \(\binom nk\) en \(P(X=k)\).
  \item Contrario de \guillemotleft{}al menos \(k\)\guillemotright{} es \guillemotleft{}como máximo \(k-1\)\guillemotright{}.
  \item Escribir la varianza como segundo parámetro de \(N(\mu,\sigma)\).
  \item Leer la cola izquierda cuando se pide la derecha.
  \item Aproximar \(P(X=k)\) por \(P(Y=k)\), que vale 0.
\end{itemize}
\end{ojo}

\end{multicols}
\end{document}
